\documentclass[a4paper,11pt]{article}
\usepackage[T2A]{fontenc}
\usepackage[utf8]{inputenc}
\usepackage[russian]{babel}
\usepackage{amsmath,amssymb,mathtools}
\usepackage{icomma}
\usepackage[a4paper,left=23mm,right=23mm,top=23mm,bottom=24mm]{geometry}
\usepackage{microtype}
\usepackage{xcolor}
\usepackage{tikz}
\usetikzlibrary{arrows.meta,positioning,shapes.geometric,calc}
\usepackage{tabularx,booktabs,array}
\usepackage{enumitem}
\usepackage{fancyhdr}
\usepackage{setspace}
\usepackage{hyperref}
\usepackage{needspace}
\definecolor{accent}{HTML}{087D83}
\definecolor{darkaccent}{HTML}{095A61}
\definecolor{textgray}{HTML}{353B3D}
\definecolor{quiet}{HTML}{687477}
\hypersetup{colorlinks=true,linkcolor=darkaccent,urlcolor=darkaccent,pdfauthor={Артём Александрович Лёвин},pdftitle={Чек-лист. Аннуитетные кредиты. Ксения Клыкова}}
\setstretch{1.17}
\setlength{\parindent}{0pt}
\setlength{\parskip}{6pt}
\setlength{\headheight}{23pt}
\setlist[enumerate]{leftmargin=*,topsep=4pt,itemsep=5pt,parsep=1pt}
\setlist[itemize]{leftmargin=*,topsep=4pt,itemsep=3pt}
\pagestyle{fancy}
\fancyhf{}
\fancyhead[L]{\footnotesize\color{quiet} Аннуитетные кредиты}
\fancyhead[R]{\footnotesize\color{quiet} Ксения Клыкова}
\fancyfoot[C]{\footnotesize\color{quiet}\thepage}
\renewcommand{\headrulewidth}{0.35pt}
\renewcommand{\footrulewidth}{0pt}
\renewcommand{\headrule}{\hbox to\headwidth{\color{accent}\leaders\hrule height \headrulewidth\hfill}}
\newcommand{\topic}[1]{\Needspace{7\baselineskip}\vspace{11pt}{\Large\color{darkaccent}#1}\par\vspace{3pt}{\color{accent}\rule{\linewidth}{0.48pt}}\vspace{6pt}}
\newcommand{\subtopic}[1]{\vspace{8pt}{\color{darkaccent}\bfseries #1}\par\vspace{2pt}}
\newcommand{\keyidea}[1]{{\color{darkaccent}\bfseries Ключевая идея.} #1}
\begin{document}
\thispagestyle{empty}
\begin{center}
\vspace*{24mm}
{\color{darkaccent}\large\MakeUppercase{Математика}}\par
\vspace{19mm}
{\Huge Чек-лист}\par
\vspace{12mm}
{\LARGE\color{accent} Аннуитетные кредиты}\par
\vspace{3mm}
{\Large\color{accent} Равные платежи и переплата}\par
\vspace{16mm}
{\normalsize Подготовка к ЕГЭ по профильной математике}\par
\vfill
{\large Ксения Клыкова}\par
\vspace{9mm}
{\normalsize 29 октября 2026 г.}\par
\vspace{17mm}
{\small Лёвин Артём Александрович}\par
{\small эксклюзивно для Ксении Клыковой}\par
\vspace{17mm}
\begin{tikzpicture}[x=1cm,y=1cm]
\draw[accent,line width=0.85pt] (0,0) .. controls (1.2,0.65) and (2.2,-0.55) .. (3.4,0)
  .. controls (4.5,0.58) and (5.65,-0.58) .. (6.9,0)
  .. controls (8.05,0.55) and (9.3,-0.55) .. (10.5,0);
\end{tikzpicture}
\vspace{9mm}
\end{center}
\clearpage
\topic{1. Как распознать аннуитетную схему}
\textbf{Сигнал в условии:} кредит погашается \textbf{равными платежами} через одинаковые промежутки времени. В задачах этого занятия платежи выполняются \textbf{в конце каждого года}, после начисления годовых процентов.

Банк начисляет проценты на \emph{остаток долга}, после чего заёмщик вносит очередной платёж. Платёж одинаков каждый год, но сумма процентов в разные годы обычно различается.

\topic{2. Обозначения и повышающий коэффициент}
\begin{tabularx}{\linewidth}{@{}lX@{}}
\toprule
\(S\) & первоначальная сумма кредита; \\
\(r\) & годовая процентная ставка в процентах; \\
\(k\) & повышающий коэффициент; \\
\(x\) & одинаковый ежегодный платёж; \\
\(B_i\) & остаток долга после платежа за \(i\)-й год; \(B_0=S\); \\
\(n\) & число лет кредитования. \\
\bottomrule
\end{tabularx}

\[
{\color{darkaccent}k=1+\frac{r}{100}},\qquad
{\color{darkaccent}B_i=kB_{i-1}-x}.
\]

\keyidea{Сначала умножаем долг на \(k\) (начисляем проценты), затем вычитаем платёж \(x\).} Если кредит погашен полностью, то \(B_n=0\).

\subtopic{Пример. Первый год}
Взяли \(200\) тыс. руб. под \(20\%\) годовых. Коэффициент \(k=1{,}2\).
\[
200\cdot0{,}2=40,\qquad 200+40=240,\qquad B_1=240-x.
\]
В начале второго года долг равен \(240-x\), поэтому проценты начисляют именно на \(240-x\), а не на первоначальные \(200\).

\topic{3. Таблица задолженности на три года}
Продолжим пример с \(S=200\), \(r=20\%\). Все суммы указаны в тысячах рублей.

\begingroup\small
\renewcommand{\arraystretch}{1.45}
\begin{tabularx}{\linewidth}{@{}c>{\centering\arraybackslash}X>{\centering\arraybackslash}X>{\centering\arraybackslash}Xc>{\centering\arraybackslash}X@{}}
\toprule
Год & Долг в начале & Проценты & Перед платежом & Платёж & Остаток \\
\midrule
1 & \(200\) & \(40\) & \(240\) & \(x\) & \(240-x\) \\
2 & \(240-x\) & \(0{,}2(240-x)\) & \(1{,}2(240-x)\) & \(x\) & \(B_2\) \\
3 & \(B_2\) & \(0{,}2B_2\) & \(1{,}2B_2\) & \(x\) & \(0\) \\
\bottomrule
\end{tabularx}
\endgroup
\[
B_2=1{,}2(240-x)-x,\qquad B_3=1{,}2B_2-x=0.
\]

\clearpage
\topic{4. Рабочие формулы для равных платежей}
\subtopic{Два года}
В конце первого года осталось \(B_1=Sk-x\). Во втором году:
\[
B_2=(Sk-x)k-x=Sk^2-xk-x=0.
\]
Соберём слагаемые с \(x\):
\[
Sk^2-x(k+1)=0,\qquad
{\color{darkaccent}x=\frac{Sk^2}{k+1}}.
\]

\subtopic{Три года}
\begin{align*}
B_1&=Sk-x,\\
B_2&=(Sk-x)k-x=Sk^2-xk-x,\\
B_3&=(Sk^2-xk-x)k-x\\
   &=Sk^3-xk^2-xk-x=0.
\end{align*}
Выносим общий множитель \(x\):
\[
Sk^3-x(k^2+k+1)=0,\qquad
{\color{darkaccent}x=\frac{Sk^3}{k^2+k+1}}.
\]

\textbf{Закономерность:} при \(n\) годах первый член содержит \(Sk^n\), а платежи вычитаются с последовательно убывающими степенями \(k\):
\[
Sk^n-x\bigl(k^{n-1}+k^{n-2}+\ldots+k+1\bigr)=0.
\]

\topic{5. Общая сумма выплат и переплата}
Если каждый год вносится \(x\), а срок составляет \(n\) лет, то
\[
\begin{aligned}
\text{Общая сумма выплат}\;&=nx,\\
\text{Переплата}\;&=nx-S,\\
\text{Процент переплаты}\;&=\frac{nx-S}{S}\cdot100\%.
\end{aligned}
\]
\keyidea{Годовая ставка \(r\) и процент переплаты за весь кредит --- разные величины.} Переплата зависит в том числе от количества платежей и срока кредита.

\topic{6. Образец расчёта платежа}
\textbf{Условие.} Кредит \(455\) тыс. руб. взят под \(20\%\) годовых на три года. Он погашается тремя равными ежегодными платежами. Найдите платёж.
\[
S=455,\quad k=1{,}2=\frac65,\quad n=3.
\]
\[
x=\frac{455\cdot1{,}2^3}{1{,}2^2+1{,}2+1}
 =\frac{455\cdot1{,}728}{3{,}64}=216\quad\text{(тыс. руб.)}.
\]
Проверка остатков: \(455\cdot1{,}2-216=330\), затем \(330\cdot1{,}2-216=180\), затем \(180\cdot1{,}2-216=0\). Долг погашен.

\topic{7. Как найти неизвестную процентную ставку}
\textbf{Условие.} Кредит \(2150\) тыс. руб. погашен двумя равными ежегодными платежами по \(1322{,}5\) тыс. руб. Найдите годовую ставку.

Для двух лет используем уравнение \(Sk^2-xk-x=0\):
\[
2150k^2-1322{,}5k-1322{,}5=0.
\]
Умножим обе части на \(2\), затем разделим на \(5\):
\[
4300k^2-2645k-2645=0
\ \Longleftrightarrow\ 860k^2-529k-529=0.
\]
Дискриминант:
\[
D=(-529)^2-4\cdot860\cdot(-529)
 =279841+1819760=2099601=1449^2.
\]
\[
k_{1,2}=\frac{529\pm1449}{1720},\qquad
k_1=\frac{1978}{1720}=1{,}15,\qquad k_2<0.
\]
Отрицательный коэффициент невозможен: при положительной годовой ставке \(k>1\). Следовательно,
\[
1+\frac{r}{100}=1{,}15\ \Longrightarrow\ {\color{darkaccent}r=15\%}.
\]

\subtopic{Как проверить большой квадрат вручную}
\[
1449^2=(1450-1)^2=1450^2-2\cdot1450+1=2099601.
\]
На экзамене полезно сначала сокращать коэффициенты уравнения, затем проверять дискриминант и допустимость найденного \(k\).

\topic{8. На что обратить внимание}
\begin{itemize}
\item \textbf{Проценты.} \(20\%=0{,}2\), значит \(k=1{,}2\). Замена на \(0{,}4\) или \(1{,}4\) меняет всю таблицу.
\item \textbf{Порядок операций.} Сначала проценты на остаток, затем вычитается \(x\).
\item \textbf{Равенство платежей.} Во всех строках один и тот же \(x\).
\item \textbf{Последний год.} После полного погашения \(B_n=0\).
\item \textbf{Корни уравнения.} Проверяйте \(k>1\) при положительной ставке и \(x>0\).
\item \textbf{Арифметика.} При дробных коэффициентах сначала избавьтесь от десятичных дробей и сократите общие множители.
\end{itemize}

\topic{9. Проверка результата по годам}
При найденной ставке \(15\%\) из примера раздела 7 можно проверить погашение без повторного решения квадратного уравнения:
\[
\begin{aligned}
B_1&=2150\cdot1{,}15-1322{,}5=1150,\\
B_2&=1150\cdot1{,}15-1322{,}5=0.
\end{aligned}
\]
Остаток равен нулю, платежи одинаковы. Значит, исходные данные согласованы с найденной ставкой.

\topic{10. Процент переплаты: короткий пример}
Если заёмщик получил \(480\) тыс. руб. и внёс три платежа по \(190\) тыс. руб., то
\[
\begin{aligned}
\text{Всего выплачено:}\quad&3\cdot190=570;\\
\text{Переплата:}\quad&570-480=90;\\
\text{Процент переплаты:}\quad&\frac{90}{480}\cdot100\%=18{,}75\%.
\end{aligned}
\]
Эти \(18{,}75\%\) описывают переплату за весь срок. Для определения \emph{годовой} ставки потребуется уравнение аннуитетной схемы.

\clearpage
\begin{center}
{\Large\color{darkaccent}Алгоритм: аннуитетный кредит}\par
\vspace{6mm}
\end{center}
\tikzset{
flowstep/.style={draw=accent,rounded corners=2pt,very thin,align=center,minimum width=10.3cm,minimum height=12mm,text width=9.5cm,fill=white,font=\normalsize},
terminal/.style={flowstep,rounded corners=7mm,minimum width=8.5cm,text width=7.8cm},
arr/.style={-{Latex[length=2.5mm]},line width=0.75pt,draw=accent}}
\begin{center}
\begin{tikzpicture}[node distance=10mm]
\node[terminal] (a) {Равные платежи \(x\) в конце каждого года};
\node[flowstep,below=of a] (b) {Запишите \(S,r,n,x\); неизвестное обозначьте буквой};
\node[flowstep,below=of b] (c) {Найдите повышающий коэффициент: \(\displaystyle k=1+\frac{r}{100}\)};
\node[flowstep,below=of c] (d) {Каждый год пересчитывайте остаток: \(\displaystyle B_i=kB_{i-1}-x\)};
\node[flowstep,below=of d] (e) {Используйте условие полного погашения: \(\displaystyle B_n=0\)};
\node[flowstep,below=of e] (f) {Раскройте скобки, вынесите \(x\) или решите уравнение относительно \(k\)};
\node[flowstep,below=of f] (g) {Проверьте \(k>1\), \(x>0\) и итоговый остаток};
\node[terminal,below=of g] (h) {Запишите ответ с единицами измерения};
\foreach \u/\v in {a/b,b/c,c/d,d/e,e/f,f/g,g/h} \draw[arr] (\u.south)--(\v.north);
\end{tikzpicture}
\end{center}

\clearpage
\topic{Тренировочный блок · Путь А}
Во всех задачах платежи вносятся в конце года после начисления процентов. Денежные суммы указаны в \textbf{тысячах рублей}, если не оговорено иное. Указывайте ключевые преобразования, ограничения и ответ.

\subtopic{Базовый уровень}
\begin{enumerate}[label=\textbf{\arabic*.},itemsep=9pt]
\item Годовая ставка кредита равна \(12\%\). Найдите повышающий коэффициент \(k\).
\item Кредит \(200\) тыс. руб. взят под \(20\%\) годовых. Найдите начисленные проценты за первый год, задолженность перед первым платежом и остаток после платежа \(x\).
\item Кредит \(300\) тыс. руб. взят под \(10\%\) годовых на два года. Запишите остатки \(B_1\), \(B_2\) через \(x\) и уравнение полного погашения.
\item Для двухлетнего кредита с первоначальной суммой \(S\) и коэффициентом \(k\) выведите формулу одинакового ежегодного платежа \(x\).
\end{enumerate}

\subtopic{Средний уровень}
\begin{enumerate}[label=\textbf{\arabic*.},start=5,itemsep=9pt]
\item Кредит \(210\) тыс. руб. взят под \(10\%\) годовых на два года с равными платежами. Найдите ежегодный платёж.
\item Кредит \(244\) тыс. руб. взят под \(25\%\) годовых на три года с равными платежами. Найдите ежегодный платёж.
\item Кредит \(600\) тыс. руб. погашен тремя равными платежами по \(250\) тыс. руб. Найдите общую сумму выплат и переплату.
\end{enumerate}

\subtopic{Уровень выше среднего}
\begin{enumerate}[label=\textbf{\arabic*.},start=8,itemsep=9pt]
\item Кредит \(100\) тыс. руб. взят под \(50\%\) годовых на два года с равными платежами. Найдите \textbf{процент переплаты} за весь срок кредита.
\item Двухлетний кредит под \(25\%\) годовых погашен двумя равными платежами по \(200\) тыс. руб. Найдите первоначальную сумму кредита.
\item Кредит \(550\) тыс. руб. погашен за два года двумя равными платежами по \(360\) тыс. руб. Найдите годовую процентную ставку. Квадратное уравнение решите без калькулятора.
\end{enumerate}

\clearpage
\topic{Ответы к тренировочному блоку}
\renewcommand{\arraystretch}{1.5}
\begin{tabularx}{\linewidth}{@{}cX@{}}
\toprule
№ & Ответ \\
\midrule
1 & \(k=1{,}12\).\\
2 & \(40\) тыс. руб.; \(240\) тыс. руб.; \(240-x\) тыс. руб.\\
3 & \(B_1=330-x\); \(B_2=1{,}1(330-x)-x=363-2{,}1x\); \(363-2{,}1x=0\).\\
4 & \(\displaystyle x=\frac{Sk^2}{k+1}\).\\
5 & \(121\) тыс. руб.\\
6 & \(125\) тыс. руб.\\
7 & Общая сумма \(750\) тыс. руб.; переплата \(150\) тыс. руб.\\
8 & \(80\%\).\\
9 & \(288\) тыс. руб.\\
10 & \(20\%\) годовых.\\
\bottomrule
\end{tabularx}

\vspace{12mm}
\textcolor{darkaccent}{\bfseries Итоговая памятка}\par\vspace{3mm}
\begin{tabularx}{\linewidth}{@{}>{\raggedright\arraybackslash}p{0.41\linewidth}X@{}}
\toprule
Сигнал & Что делать \\
\midrule
Равные ежегодные платежи & Используйте аннуитетную схему. \\
Ставка \(r\%\) & Замените на \(k=1+r/100\). \\
Остаток после года & \(B_i=kB_{i-1}-x\). \\
Полное погашение & Приравняйте \(B_n\) к нулю. \\
Два года & \(Sk^2-xk-x=0\). \\
Три года & \(Sk^3-xk^2-xk-x=0\). \\
Общая сумма выплат & \(nx\). \\
Переплата & \(nx-S\). \\
\bottomrule
\end{tabularx}

\vfill
\begin{center}
{\small\color{quiet}Лёвин Артём Александрович эксклюзивно для Ксении Клыковой}
\end{center}
\end{document}
